\exercisesection{\S~2}

\begin{enumerate}
\def\labelenumi{\arabic{enumi})}
\tightlist
\item
  设 \(\mathbf{G}\) 是一个连通紧李群，\(g\) 是 \(\mathbf{G}\) 的一个元素。
\end{enumerate}

\begin{enumerate}
\def\labelenumi{\alph{enumi})}
\item
  证明存在一个整数 \(n \geq 1\)，使得中心化子 \(Z(g^{n})\) 是连通的（证明对于适当的 n，由 \(g^{n}\) 生成的闭子群是一个环面）。
\item
  假设 \(\operatorname{Ker}(\operatorname{Ad}g^n - 1)\) 的维数与 \(n\) 无关（\(n \geq 1\)）。证明 \(Z(g)\) 是连通的。
\item
  如果 \(g^n\) 对于所有 \(n \geq 1\) 都是正则元，则 \(Z(g)\) 是连通的。
\end{enumerate}

\begin{enumerate}
\def\labelenumi{\arabic{enumi})}
\setcounter{enumi}{1}
\item
  证明每个连通紧李群 G 都是其导群与一个环面的半直积（注意，若 T 是 G 的极大环面，则 \(D(G) \cap T\) 是一个环面）。
\item
  设 G 是一个紧李群，\(\mathfrak{g}\) 是其李代数，\(\mathfrak{s}\) 是 \(\mathfrak{g}\) 的向量子空间，满足对于 \(x, y, z\)（它们属于 \(\mathfrak{s}\)），有 \([x, [y, z]] \in \mathfrak{s}\)。令 \(\mathcal{L}\) 为 \(\mathfrak{g}\) 中包含于 \(\mathfrak{s}\) 的交换子代数集合。证明 G 中 \(\mathfrak{s}\) 的稳定化子的单位元连通分支在 \(\mathcal{L}\) 的极大元素集合上传递地作用（论证如定理 1 的证明）。
\item
  设 G 是一个连通紧李群，T 是 G 的极大环面，S 是 T 的子环面。分别记 \(\Sigma\)（分别为 F）为 S 在 \(W_{G}(T)\) 中的稳定化子（分别为固定子）。
\end{enumerate}

\begin{enumerate}
\def\labelenumi{\alph{enumi})}
\item
  证明群 \(\mathrm{N_G(S) / Z_G(S)}\) 同构于商 \(\Sigma /\mathrm{F}\)。
\item
  设 H 是 G 的连通闭子群，且 S 是 H 的极大环面。证明 \(W_{\mathrm{H}}(\mathrm{S})\) 的每个元素，视为 S 的自同构，都是 \(\Sigma\) 中某个元素在 S 上的限制。
\end{enumerate}

\begin{enumerate}
\def\labelenumi{\arabic{enumi})}
\setcounter{enumi}{4}
\tightlist
\item
  设 \(G\) 是一个连通紧李群，且 \(S\) 是 \(G\) 的一个环面。证明下列条件等价：
\end{enumerate}

\begin{enumerate}
\def\labelenumi{(\roman{enumi})}
\item
  S 包含于唯一的极大环面中；
\item
  \(Z_{G}(S)\) 是一个极大环面；
\item
  S 包含一个正则元。
\end{enumerate}

\begin{enumerate}
\def\labelenumi{\arabic{enumi})}
\setcounter{enumi}{5}
\tightlist
\item
  设 G 是一个连通紧李群，T 是 G 的极大环面，\(\mathfrak{g}\)（分别为 \(\mathfrak{t}\)）是 G（分别为 T）的李代数，且 \(i : \mathfrak{t} \to \mathfrak{g}\) 是典则嵌入。证明映射 \(^{t}i : \mathfrak{g}^{*} \to \mathfrak{t}^{*}\) 通过取商诱导出从 \(\mathfrak{g}^{*}/G\) 到 \(\mathfrak{t}^{*}/W_{G}(T)\) 的一个同胚。
\end{enumerate}

¶ 7) 设 X 和 Y 是两个类 \(C^{r}\) 的实流形（\(1 \leq r \leq \omega\)），它们是分离的、连通的且维数为 n；设 \(f : X \to Y\) 是一个类 \(C^{r}\) 的真态射，并配备定向 F（《微分流形与解析流形》，结果 10.2.5）。

\begin{enumerate}
\def\labelenumi{\alph{enumi})}
\item
  证明存在唯一实数 d，使得 \(\int_{X} f^{*}\alpha = d\int_{Y}\alpha\) 对 Y 上每个扭曲微分形式 \(\alpha\)（次数为 n、类 \(C^{1}\) 且具有紧支集）成立（使用《微分流形与解析流形》，结果 11.2.4）。
\item
  设 \(y \in \mathbf{Y}\) 满足 \(f\) 在 \(f^{-1}(y)\) 的每一点处都无分歧。对于 \(x \in f^{-1}(y)\)，令 \(\nu_x(f) = 1\)（分别令 \(\nu_x(f) = -1\)），若映射 \(\mathrm{F}_x\) 和 \(\tilde{f}_x\)（《微分流形与解析流形》，结果 10.2.5，例 b）从 \(\mathrm{Or}(\mathrm{T}_x(\mathrm{X}))\) 到 \(\mathrm{Or}(\mathrm{T}_y(\mathrm{Y}))\) 相同（分别为相反映射）。证明 \(d = \sum_{x \in f^{-1}(y)} \nu_x(f)\)，特别地，\(d \in \mathbf{Z}\)。
\end{enumerate}

称 d 为 f 的次数，并记为 \(\deg f\)。若 X = Y 且 X 可定向，则可以取 F 为保持 X 的定向的定向。

\begin{enumerate}
\def\labelenumi{\alph{enumi})}
\setcounter{enumi}{2}
\item
  证明，如果 \(\deg f \neq 0\)，则 \(f\) 是满射。
\item
  如果存在 \(x \in X\)，使得 \(f^{-1}(f(x)) = \{x\}\)，且 f 在 x 处无分歧，则 f 是满射。
\end{enumerate}

\begin{enumerate}
\def\labelenumi{\arabic{enumi})}
\setcounter{enumi}{7}
\tightlist
\item
  设 G 是一个连通紧李群，dg 是 G 上的归一化 Haar 测度。
\end{enumerate}

\begin{enumerate}
\def\labelenumi{\alph{enumi})}
\tightlist
\item
  令 \(f: \mathrm{G} \to \mathrm{G}\) 为一个流形态射；对于 \(g \in \mathrm{G}\)，微分 \(\mathrm{T}_g(f)\) 可通过左平移识别为从 \(\mathrm{L}(\mathrm{G})\) 到 \(\mathrm{L}(\mathrm{G})\) 的线性映射。证明下列公式
\end{enumerate}

\[
\deg f. \int_ {\mathrm{G}} \varphi (g) d g = \int_ {\mathrm{G}} \varphi (f (g)) \det \mathrm{T} _ {g} (f) d g
\]

对 G 上每个可积的（复值）函数 \(\varphi\)，以及

\[
\deg f = \int_ {\mathrm{G}} \det \mathrm{T} _ {g} (f) d g.
\]

\begin{enumerate}
\def\labelenumi{\alph{enumi})}
\setcounter{enumi}{1}
\tightlist
\item
  令 \(\psi_k: \mathbf{G} \to \mathbf{G}\) 为映射 \(g \mapsto g^k\)。证明公式
\end{enumerate}

\[
\deg \psi_ {k} = \int_ {\mathrm{G}} \det (1 + \operatorname{Ad} g + \dots + (\operatorname{Ad} g) ^ {k - 1}) d g.
\]

\begin{enumerate}
\def\labelenumi{\alph{enumi})}
\setcounter{enumi}{2}
\tightlist
\item
  令 \(\mathrm{T}\) 是 \(\mathrm{G}\) 的极大环面，且 \(\mathrm{T}_r\) 是 \(\mathrm{T}\) 的正则元集合。证明 \(\exp_{\mathrm{G}}^{-1}(\mathrm{T}_r) \subset \mathrm{L}(\mathrm{T})\)，并且 \(\psi_k^{-1}(\mathrm{T}_r) \subset \mathrm{T}_r\) 对所有 \(k \geq 1\) 成立。
\end{enumerate}

\(d\)) 证明 \(\deg \psi_{k} = k^{\dim (\mathrm{T})}\)；推出等式

\[
\int_ {\mathrm{G}} \det (1 + \operatorname{Ad} g + \dots + (\operatorname{Ad} g) ^ {k - 1}) d g = k ^ {\dim (\mathrm{T})}.
\]

\begin{enumerate}
\def\labelenumi{\arabic{enumi})}
\setcounter{enumi}{8}
\tightlist
\item
  本习题旨在给出定理 2 的另一证明。设 G 是一个连通紧李群，T 是 G 的极大环面，N 是 T 的正规化子。记 \(G \times^{N} T\) 为 \(G \times T\) 按 N 的作用 \((g, t).n = (gn, n^{-1}tn)\) 得到的商流形（《微分流形与解析流形》，结果 6.5.1）。
\end{enumerate}

\begin{enumerate}
\def\labelenumi{\alph{enumi})}
\item
  证明态射 \((g, t) \mapsto gtg^{-1}\) 从 \(G \times T\) 到 G，通过取商定义出一个解析态射 \(f : G \times^{N} T \to G\)。
\item
  设 \(\theta\) 是 T 中一个元素，其幂集在 T 中稠密（《一般拓扑》，第 VII 章，§1，节 3，命题 7 的推论 2）。证明 \(f^{-1}(\theta) = \{x\}\)，其中 \(x\) 是 \((e,\theta)\) 在 \(\mathrm{G} \times^{\mathrm{N}}\mathrm{T}\) 中的类，并且 \(f\) 在 \(x\) 处无分歧。
\item
  由习题 7 d) 得出 \(f\) 是满射，并由此给出定理 2 的另一证明。
\end{enumerate}

\begin{enumerate}
\def\labelenumi{\arabic{enumi})}
\setcounter{enumi}{9}
\tightlist
\item
  * 本习题使用代数拓扑中的下列结果（Lefschetz 公式 \(^{11}\)）：设 X 是有限维紧流形，f 是从 X 到自身的一个态射。假设 X 中满足 \(f(x) = x\) 的点 x 所成的集合 F 是有限的，并且对于所有 \(x \in F\)，数 \(\delta(x) = \det(1 - \mathrm{T}_{x}(f))\) 都非零。如果 \(\mathrm{H}^{i}(f)\) 表示由 f 诱导的向量空间 \(\mathrm{H}^{i}(\mathrm{X}, \mathbf{R})\) 的自同态（对于 \(i \geq 0\)），则
\end{enumerate}

\[
\sum_ {x \in \mathrm{F}} \delta (x) / | \delta (x) | = \sum_ {i \geq 0} (- 1) ^ {i} \operatorname{Tr} \mathrm{H} ^ {i} (f).
\]

设 G 是一个连通紧李群，T 是 G 的极大环面。对于 \(g \in G\)，记 \(\tau(g)\) 为由 g 的左乘法在流形 G/T 上诱导的自同构。

\begin{enumerate}
\def\labelenumi{\alph{enumi})}
\item
  设 \(t\) 是 \(\mathrm{T}\) 中的一个元素，使得由 \(t\) 生成的子群在 \(\mathrm{T}\) 中稠密。证明 \(\tau(t)\) 的不动点正是 \(n\mathrm{T}\)（其中 \(n \in \mathrm{N}_{\mathrm{G}}(\mathrm{T})\)）的类。从而 \((\mathrm{N}_{\mathrm{G}}(\mathrm{T}):\mathrm{T}) = \sum_{i}(-1)^{i}\dim_{\mathbf{R}}\mathrm{H}^{i}(\mathrm{G / T},\mathbf{R})\)。
\item
  令 g 是 G 的任意元素；证明 \(\tau(g)\) 的不动点集非空，并由此给出定理 2 的另一证明。*
\end{enumerate}

\begin{enumerate}
\def\labelenumi{\arabic{enumi})}
\setcounter{enumi}{10}
\tightlist
\item
  设 G 是一个紧李群。如果 G 的子群 S 等于某个在其正规化子中具有有限指数的循环子群的闭包，则称 S 为 (C) 型子群。
\end{enumerate}

\begin{enumerate}
\def\labelenumi{\alph{enumi})}
\item
  设 \(S\) 是一个 (C) 型子群；证明 \(S_0\) 是 \(G_0\) 的极大环面，且 \(S\) 是 \(S_0\) 与某个有限循环子群的直积。
\item
  证明 G 的每个元素 g 都包含于一个 (C) 型子群中（考虑由 g 和 \(Z(g)_{0}\) 的一个极大环面生成的群）。
\item
  设 S 是一个 (C) 型子群，且 s 是 S 中幂集在 S 中稠密的元素。证明 \(sG_{0}\) 中的每个元素都在 \(\operatorname{Int}(G_{0})\) 下与 \(sS_{0}\) 中的某个元素共轭（使用习题 10 的方法）。
\item
  记 \(p: G \to G/G_{0}\) 为商映射。证明映射 \(S \mapsto p(S)\) 从 G 的 (C) 型子群的共轭类集合到 \(G/G_{0}\) 的循环子群的共轭类集合诱导出一个双射。
\item
  设 S 是 G 的一个 (C) 型子群；记 \(S_{\rho}\) 为 S 中其在 \(S/S_{0}\) 中的类生成 \(S/S_{0}\) 的元素集合。证明 \(S_{\rho}\) 中在 G 中共轭的两个元素在 \(N_{G}(S)\) 中共轭。
\end{enumerate}
% semantic-fragments: legacy-input-seam
\relax{}
